\documentclass{amsart}
% For dealing with references we use the comment environment
\usepackage{verbatim}
\newenvironment{reference}{\comment}{\endcomment}
%\newenvironment{reference}{}{}
\newenvironment{slogan}{\comment}{\endcomment}
\newenvironment{history}{\comment}{\endcomment}

% For commutative diagrams we use Xy-pic
\usepackage[all]{xy}

% We use 2cell for 2-commutative diagrams.
\xyoption{2cell}
\UseAllTwocells

% We use multicol for the list of chapters between chapters
\usepackage{multicol}

% This is generally recommended for better output
\usepackage{lmodern}
\usepackage[T1]{fontenc}

% For cross-file-references

% Package for hypertext links:
\usepackage{hyperref}

% For any local file, say "hello.tex" you want to link to please

% Theorem environments.
%
\theoremstyle{plain}
\newtheorem{theorem}[subsection]{Theorem}
\newtheorem{proposition}[subsection]{Proposition}
\newtheorem{lemma}[subsection]{Lemma}

\theoremstyle{definition}
\newtheorem{definition}[subsection]{Definition}
\newtheorem{example}[subsection]{Example}
\newtheorem{exercise}[subsection]{Exercise}
\newtheorem{situation}[subsection]{Situation}

\theoremstyle{remark}
\newtheorem{remark}[subsection]{Remark}
\newtheorem{remarks}[subsection]{Remarks}

\numberwithin{equation}{subsection}

% Macros
%
\def\lim{\mathop{\mathrm{lim}}\nolimits}
\def\colim{\mathop{\mathrm{colim}}\nolimits}
\def\Spec{\mathop{\mathrm{Spec}}}
\def\Hom{\mathop{\mathrm{Hom}}\nolimits}
\def\Ext{\mathop{\mathrm{Ext}}\nolimits}
\def\SheafHom{\mathop{\mathcal{H}\!\mathit{om}}\nolimits}
\def\SheafExt{\mathop{\mathcal{E}\!\mathit{xt}}\nolimits}
\def\Sch{\mathit{Sch}}
\def\Mor{\mathop{\mathrm{Mor}}\nolimits}
\def\Ob{\mathop{\mathrm{Ob}}\nolimits}
\def\Sh{\mathop{\mathit{Sh}}\nolimits}
\def\NL{\mathop{N\!L}\nolimits}
\def\CH{\mathop{\mathrm{CH}}\nolimits}
\def\proetale{{pro\text{-}\acute{e}tale}}
\def\etale{{\acute{e}tale}}
\def\QCoh{\mathit{QCoh}}
\def\Ker{\mathop{\mathrm{Ker}}}
\def\Im{\mathop{\mathrm{Im}}}
\def\Coker{\mathop{\mathrm{Coker}}}
\def\Coim{\mathop{\mathrm{Coim}}}

% Boxtimes
%
\DeclareMathSymbol{\boxtimes}{\mathbin}{AMSa}{"02}

%
% Macros for moduli stacks/spaces
%
\def\QCohstack{\mathcal{QC}\!\mathit{oh}}
\def\Cohstack{\mathcal{C}\!\mathit{oh}}
\def\Spacesstack{\mathcal{S}\!\mathit{paces}}
\def\Quotfunctor{\mathrm{Quot}}
\def\Hilbfunctor{\mathrm{Hilb}}
\def\Curvesstack{\mathcal{C}\!\mathit{urves}}
\def\Polarizedstack{\mathcal{P}\!\mathit{olarized}}
\def\Complexesstack{\mathcal{C}\!\mathit{omplexes}}
% \Pic is the operator that assigns to X its picard group, usage \Pic(X)
% \Picardstack_{X/B} denotes the Picard stack of X over B
% \Picardfunctor_{X/B} denotes the Picard functor of X over B
\def\Pic{\mathop{\mathrm{Pic}}\nolimits}
\def\Picardstack{\mathcal{P}\!\mathit{ic}}
\def\Picardfunctor{\mathrm{Pic}}
\def\Deformationcategory{\mathcal{D}\!\mathit{ef}}

\setlength{\emergencystretch}{3em}
\begin{document}
\title[Stacks Project, Tag 0C0G]{249 --- Identification of upper shriek with derived Hom for finite ring maps}
\maketitle
\noindent Copyright (C) 2005--2025 Johan de Jong. Stacks Project authors. Source: \href{https://stacks.math.columbia.edu/tag/0C0G}{Tag 0C0G}.

\noindent Editorial difficulty note (not author text). Difficulty H2: The proof reduces a general finite algebra to a monogenic quotient by a monic polynomial using transitivity. It constructs the dual basis action explicitly and shows the coefficient functional generates the algebra dual as a rank-one module. The resulting comparison identifies independently constructed derived functors, not merely their underlying complexes..

\noindent Editorial provenance note (not author text). History: excerpt from revision \texttt{a0444\allowbreak 6e57e\allowbreak c1fbc\allowbreak 252a8\allowbreak 71afc\allowbreak ec775\allowbreak 2fb28\allowbreak 07b14}, dualizing.tex, lines 4598--4645. Reference presentation was adapted; original-author context and core support blocks were retained or added. The mathematical wording is unchanged. Licensed under GNU FDL 1.2 or later; no invariant sections or cover texts. See ../licenses/COPYING. Context blocks reproduce author definitions and supporting results; they are not additional counted problems.

\section{Upper shriek algebraically}
\label{section-relative-dualizing-complex-algebraic}

\noindent
For a finite type homomorphism $R \to A$ of Noetherian rings
we will construct a functor $\varphi^! : D(R) \to D(A)$
well defined up to nonunique isomorphism which
as we will see in Duality for Schemes, Remark
\href{https://stacks.math.columbia.edu/tag/0BV2}{remark local calculation shriek (Tag 0BV2)}
agrees up to isomorphism with the upper shriek functors
one encounters in the duality theory for schemes.
To motivate the construction we mention two additional properties:
\begin{enumerate}
\item $\varphi^!$ sends a dualizing complex for $R$ (if it exists)
to a dualizing complex for $A$, and
\item $\omega_{A/R}^\bullet = \varphi^!(R)$ is a kind of
relative dualizing complex: it lies in $D^b_{\textit{Coh}}(A)$ and restricts
to a dualizing complex on the fibres provided $R \to A$ is flat.
\end{enumerate}
These statements are Lemmas \href{https://stacks.math.columbia.edu/tag/0BZL}{lemma shriek dualizing algebraic (Tag 0BZL)} and
\href{https://stacks.math.columbia.edu/tag/0BZW}{lemma relative dualizing algebraic (Tag 0BZW)}.

\medskip\noindent
Let $\varphi : R \to A$ be a finite type homomorphism of Noetherian rings.
We will define a functor $\varphi^! : D(R) \to D(A)$ in the following way
\begin{enumerate}
\item If $\varphi : R \to A$ is surjective we set
$\varphi^!(K) = R\Hom(A, K)$. Here we use the functor
$R\Hom(A, -) : D(R) \to D(A)$ of
Section \href{https://stacks.math.columbia.edu/tag/0A6Z}{section trivial (Tag 0A6Z)}, and
\item in general we choose a surjection $\psi : P \to A$ with
$P = R[x_1, \ldots, x_n]$ and we set
$\varphi^!(K) = \psi^!(K \otimes_R^\mathbf{L} P)[n]$.
Here we use the functor
$- \otimes_R^\mathbf{L} P : D(R) \to D(P)$
of More on Algebra, Section \href{https://stacks.math.columbia.edu/tag/06Y5}{section derived base change (Tag 06Y5)}.
\end{enumerate}
Note the shift $[n]$ by the number of variables in the polynomial
ring. This construction is {\bf not} canonical and the functor
$\varphi^!$ will only be well defined up to a (nonunique) isomorphism of
functors\footnote{It is possible to make the construction canonical:
use $\Omega^n_{P/R}[n]$ instead of $P[n]$ in the
construction and use this in Lemma \href{https://stacks.math.columbia.edu/tag/0BZJ}{lemma well defined (Tag 0BZJ)}.
The material in this section becomes a lot more involved
if one wants to do this.}.



\begin{lemma}
\label{lemma-upper-shriek-finite}
Let $\varphi : R \to A$ be a finite map of Noetherian rings.
Then $\varphi^!$ is isomorphic to the functor
$R\Hom(A, -) : D(R) \to D(A)$ from
Section \href{https://stacks.math.columbia.edu/tag/0A6Z}{section trivial (Tag 0A6Z)}.
\end{lemma}

\begin{proof}
Suppose that $A$ is generated by $n > 1$ elements over $R$.
Then can factor $R \to A$ as a composition of two finite ring maps
where in both steps the number of generators is $< n$.
Since we have Lemma \href{https://stacks.math.columbia.edu/tag/0BZT}{lemma composition shriek algebraic (Tag 0BZT)} and
Lemma \href{https://stacks.math.columbia.edu/tag/0C0F}{lemma composition right adjoints (Tag 0C0F)}
we conclude that it suffices
to prove the lemma when $A$ is generated by one element over $R$.
Since $A$ is finite over $R$, it follows that $A$ is a quotient
of $B = R[x]/(f)$ where $f$ is a monic polynomial in $x$
(Algebra, Lemma \href{https://stacks.math.columbia.edu/tag/00GK}{lemma finite is integral (Tag 00GK)}).
Again using the lemmas on composition and the fact that we
have agreement for surjections by definition, we conclude that
it suffices to prove the lemma for $R \to B = R[x]/(f)$.
In this case, the functor $\varphi^!$ is isomorphic to
$K \mapsto K \otimes_R^\mathbf{L} B$; you prove this by
using Lemma \href{https://stacks.math.columbia.edu/tag/0BZH}{lemma compute for effective Cartier algebraic (Tag 0BZH)}
for the map $R[x] \to B$ (note that the shift in the definition
of $\varphi^!$ and in the lemma add up to zero).
For the functor $R\Hom(B, -) : D(R) \to D(B)$ we can use
Lemma \href{https://stacks.math.columbia.edu/tag/0BZE}{lemma RHom is tensor special (Tag 0BZE)}
to see that it suffices to show $\Hom_R(B, R) \cong B$
as $B$-modules. Suppose that $f$ has degree $d$.
Then an $R$-basis for $B$ is given by $1, x, \ldots, x^{d - 1}$.
Let $\delta_i : B \to R$, $i = 0, \ldots, d - 1$
be the $R$-linear map which picks off the coefficient
of $x^i$ with respect to the given basis. Then
$\delta_0, \ldots, \delta_{d - 1}$ is a basis for $\Hom_R(B, R)$.
Finally, for $0 \leq i \leq d - 1$ a computation shows that
$$
x^i \delta_{d - 1} =
\delta_{d - 1 - i} + b_1 \delta_{d - i} + \ldots + b_i \delta_{d - 1}
$$
for some $c_1, \ldots, c_d \in R$\footnote{If
$f = x^d + a_1 x^{d - 1} + \ldots + a_d$, then
$c_1 = -a_1$, $c_2 = a_1^2 - a_2$, $c_3 = -a_1^3 + 2a_1a_2 -a_3$, etc.}.
Hence $\Hom_R(B, R)$ is a principal $B$-module with generator
$\delta_{d - 1}$. By looking
at ranks we conclude that it is a rank $1$ free $B$-module.
\end{proof}
\end{document}
